% !TeX root = ../papers/nc-p3.tex
\section{滑らかな二次曲面の各閉軌道への全退化}
\label{sec:quadric-attraction-fibers}
\begin{theorem}\label{thm:quadric-attraction-fibers}
前節の二次曲面窓を $W(m)$ とする。その軌道閉包への逆像、すなわち
$W(m)$ を軌道閉包に含む全窓の添字保持同型類は、
$m^2\ne-1$ では五つ、$m^2=-1$ では三つである。全て標準Koszul左右AS正則である。
次の表示は背景の点対応を保つという仮定を置かない。
$b$ を $X,Y,Z,W$ の一文字とし、他の平面の二文字を固定順に $c,d$ として
\[
 E_b=bcbd-bdbc-cbdb+dbcb
\]
と書く。四文字語はここでは四テンソルである。$W(m)$ を前節のClifford最高traceで規格化する。
全代表は $W(m)+A E_b+B E_{b'}$、$b\in\{X,Y\}$, $b'\in\{Z,W\}$、
$A,B\in\{0,1\}$ から得られる。残る同型は非零欄の台に対する両平面の同時反転、
および $m^2=-1$ の場合の四巡回である。
一般の場合、空台と二つの単独欄の同型類の最小周期は1、二つの両欄の類は2である。
例外 $m^2=-1$ では全三類が最小周期1である。同じファイバー内の全退化は、
これらの置換の後の台の包含である。
\end{theorem}
\begin{proof}[全パラメータでの横断接]
各文字に $(1,0),(-1,0),(0,1),(0,-1)$ を与える。
前節の軌道安定化群の連結部分はこの二次元トーラスである。実際、二つの非退化二平面を
保つ共通相似倍率の直交変換は、射影倍率を除けばこのトーラスと有限個の入替えとなる。
二平面を交換するものは $m^2=-1$ の場合にだけあり、片方の文字順を逆にするので
四文字上の四巡回となる。基点と各 $E_b$ はこの有限群とトーラスで張る部分空間内に移る。

\path{computations/verify_quadric_attraction.py} に接次元上界と持上げの有理式検算を置く。
ここで有限体での階数を全閉点へ外挿しないため、計算する小行列を記す。
各切断で $4i+j$ 順の六行六列 $P=(1,2,3,4,6,7)$ の行列式は
$-4m^2(m-1)^2(m+1)^2$ であり、許容パラメータ全域で単元である。
各切断の背景行列を $M$, $I=P^c$, $J=M_{PP}^{-1}$ とし、
$L$ を $I$ 行選択から $M_{IP}J$ 倍の $P$ 行選択を引いた行列、
$R$ を $I$ 列選択から $P$ 列選択の $JM_{PI}$ 倍を引いた行列とする。
テンソル変位の切断を $Y$ と書くと一次式は $LYR$ である。

二切断の一次式を語の指標で分け、さらに四位置に全行列単位を掛けた軌道接を引くと、
$(\pm2,0),(0,\pm2)$ には各高々一次元、零指標には高々二次元、
残りの全非零指標には何もない。具体的には各ブロックで $m=2$ での簡約階段形から
独立な列を選び、その列の転置から独立な行を選ぶ。一次式と軌道接の双方について、
選んだ小行列式を $\mathbb Q(m)$ で因数分解すると全て $m,m-1,m+1$ だけの
整数冪と非零定数の積である。従って上界は全パラメータに有効である。
四つの $E_b$ は各々一次式を消し、対応する軌道接に一列として加えた単元小行列により
横断四方向を実現する。
\end{proof}

\begin{proof}[四平面と全軌道]
余指標で基点へ流れるには四指標の各対から高々一方を選べる。
この二方向の非零台はそれぞれの文字を正にする $(\pm1,\pm1)$ で実現できる。
これらの表で変位の余重みは2、全方程式の最大余重みは4である。
実際、上の主行列式は基点で非零なので指標零、縁付き行列式の指標は余分な行と列の
四文字の指標である。重み4のSchur補正は $-LY_{*,P}JY_{P,*}R$ であり、
$Y$ を $E_b$ 一つ、および異なる平面から一つずつ選んだ二つの和にすると零になる。
前掲コードは二つの純平方項と偏極した混合項の全成分を記号的に消去する。
従って四つの表示平面は追加の非線型項を要さず、両隣接階数は正確に6である。

安定化群不変な線型横断面にこれら四方向を含めてLunaの切片を作る。
定理\ref{thm:quantum-face-resonance}の証明と同様に、完備局所環は四つの線型射影と
零指標の生成元で位相的に生成される。切片で基点へ流れる任意の点は、同じ余指標の弧上で
零指標が背景値、四射影が表示平面の点と同じ単項式であるから、その平面の点と一致する。
飽和開での強エタール性は閉軌道への全ファイバーを覆う。
各平面からの一径数退化に命題\ref{prop:unrestricted-good-reduction}を適用すれば全次数で正則である。

二平面から一つずつ選んだ指標は独立なので、代数閉体上では非零係数を全て1にできる。
$m^2\ne-1$ の場合には同時反転だけで割る。台の数は空台が1、単独が2、両方が2となる。
$m^2=-1$ の四巡回は単独四台および両方四台のそれぞれで推移的である。
一位置巡回で背景を前の表示に戻す変換は一平面内の入替えである（関係の $m$ が $-m$ に変わる）。
従って一つの指標成分だけを逆転する。これは一般の場合の両方の台だけを交換する。
形式的にもこの作用を確かめるには巡回を加えた有限拡大で切片と接射影を選べばよい。
四線の非零欄に掛かる倍率はトーラスで吸収できる。一般の場合の両台は二移動で戻り、
例外では一移動ですべて元の類に戻る。台の包含による退化は、各表示平面の二つの
独立な指標のトーラス軌道閉包から従う。
\end{proof}
